\chapter*{Úvod}\label{chap:uvod}
\addcontentsline{toc}{chapter}{Úvod}

Flash tanks are widely used in industrial systems as separation vessels designed to divide a two-phase flow into its vapor and liquid components \cite{separation_general}. Their operation is based on a sudden pressure reduction of a liquid near saturation conditions, leading to flash evaporation and the formation of a vapor–liquid mixture \cite{flash_evaporation}. In addition to flashing condensate streams, flash tanks may also receive high-pressure steam directly, for example during bypass operation, further contributing to the internal two-phase flow \cite{epk_design}. The vessel provides sufficient volume and appropriate flow conditions to enable phase separation, typically driven by gravity and density differences \cite{separation_general,flash_evaporation}. In steam power systems, this principle is applied to condensate and steam streams discharged from high-pressure regions to lower-pressure environments \cite{epk_design}. In this context, flash tanks (EPK) ensure separation and proper routing of both phases, preventing two-phase flow from entering downstream equipment such as condensers or pumps \cite{epk_design}.

Despite their apparent simplicity, flash tank operation involves complex multiphase flow phenomena governed by inertia, gravity, turbulence, and phase interaction, all strongly influenced by vessel geometry and inlet configuration \cite{multiphase_fundamentals}. In practice, design is based on empirical rules and simplified assumptions, such as velocity limits or residence time \cite{design_guidelines}. While robust, these approaches do not capture detailed internal flow structures, including recirculation, phase entrainment, or localized high-velocity regions \cite{multiphase_fundamentals,design_guidelines}.

This limitation is increasingly important in modern power systems, where higher efficiency and flexible operation impose stricter performance requirements \cite{modern_power_systems}. Off-design conditions, such as transient loads or varying inlet parameters, can significantly affect flow behavior and reduce separation efficiency, potentially leading to steam carryover or liquid entrainment and consequent performance degradation or equipment damage \cite{modern_power_systems}.

The main objective of this thesis is to analyze the internal flow behavior in a flash tank (EPK) using CFD simulations, with focus on the influence of geometric configuration and operating conditions on phase separation efficiency. Based on the results, commonly used design assumptions are assessed and potential design improvements are proposed. The work contributes to a better understanding of multiphase flow behavior in flash tanks and supports the development of more efficient and reliable designs.